\documentclass[hidelinks]{article}
\usepackage{stex}

\libinput{preambleCS}

\title{Deterministic PDAs and CFL Membership}
\author{}
\date{}

\begin{document}

\maketitle

\begin{smodule}[title={Deterministic PDAs and CFL Membership}]{DeterministicPDAsAndCFLMembership}
\usemodule{ComputerScience/FormalLanguagesAndAutomata/FormalLanguagesAndAutomata_defs?FormalLanguagesAndAutomata}
\usemodule{ComputerScience/LexingParsingCompilingInterpreting/CompilationPhases/CompilationPhases_defs?Compilation}
\usemodule{ComputerScience/LexingParsingCompilingInterpreting/Grammars/Grammars_defs?Grammars}
\usemodule{ComputerScience/LexingParsingCompilingInterpreting/Parsing/Parsing_defs?Parsing}

\section{Overview}

These notes rewrite Theory 2 Lecture 17 into prose while keeping the main technical details from the slides.
The lecture has two main parts:
\begin{itemize}
    \item \sr{deterministic pushdown automaton}{deterministic pushdown automata} (DPAs); and
    \item the \sr{CYK membership algorithm}{CYK algorithm} for testing membership in a \sn{context-free language}.
\end{itemize}

The first part studies how determinism changes the power of \sr{pushdown automaton}{pushdown automata}.
The second part gives a dynamic-programming algorithm for deciding whether a given string belongs to the language of a \sr{Compilation?ContextFreeGrammar}{context-free grammar} in \sn{Chomsky normal form}.

\section{Deterministic Pushdown Automata}

\subsection{Why deterministic PDAs matter}

The \sr{pushdown automaton}{pushdown automata} considered earlier in the course were \sr{nondeterministic pushdown automaton}{nondeterministic pushdown automata}.
That is the right level of generality for all context-free languages, but deterministic pushdown automata are still important.
In particular, parsers for programming-language syntax behave like DPAs.
For that reason, DPAs give insight into which language constructs are suitable for deterministic parsing and which ones inherently require nondeterminism.

\subsection{What determinism means for a PDA}

A \sn{deterministic pushdown automaton} is a \sn{pushdown automaton} for which there is never a genuine choice of move.
The slides formulate this in two parts:
\begin{enumerate}
    \item for a given state, input symbol (including $\lambda$), and stack symbol, there is at most one possible transition; and
    \item there is never a choice between consuming a real input symbol and taking a competing $\lambda$-transition.
\end{enumerate}

Unlike deterministic finite-state automata, DPAs are still allowed to have $\lambda$-transitions.
The point is not to ban them, but to ensure that they do not introduce branching choices.

\subsection{Formal definition}

A \sn{deterministic pushdown automaton} has the same underlying form as a \sn{nondeterministic pushdown automaton}:
\[
P = \langle Q, \Sigma, \Gamma, \delta, q_0, z_0, F \rangle
\]
with transition function
\[
\delta : Q \times (\Sigma \cup \{\lambda\}) \times \Gamma \to 2^{Q \times \Gamma^*}.
\]

The machine is deterministic if and only if the following two conditions hold:
\begin{enumerate}
    \item $\delta(q,a,X)$ has at most one member for every $q \in Q$, $a \in \Sigma \cup \{\lambda\}$, and $X \in \Gamma$; and
    \item if $\delta(q,a,X)$ is non-empty for some $a \in \Sigma$ and $X \in \Gamma$, then $\delta(q,\lambda,X)$ must be empty.
\end{enumerate}

\subsection{A DPA example}

The lecture gives a DPA for the language
\[
L = \{wcw^R \mid w \in \{0,1\}^*\},
\]
where $w^R$ denotes string reversal.

Its transitions can be grouped as follows.

\paragraph{Phase 1: read the first half and push it onto the stack.}
At state $q_0$ the machine reads symbols before the middle marker $c$ and pushes them:
\begin{itemize}
    \item $0,z_0/0z_0$
    \item $0,1/01$
    \item $0,0/00$
    \item $1,z_0/1z_0$
    \item $1,1/11$
    \item $1,0/10$
\end{itemize}

\paragraph{Phase 2: detect the middle.}
The symbol $c$ marks the split between the two halves and moves the machine from $q_0$ to $q_1$ without changing the stack content:
\begin{itemize}
    \item $c,1/1$
    \item $c,0/0$
    \item $c,z_0/z_0$
\end{itemize}

\paragraph{Phase 3: match the second half against the stack.}
At $q_1$ the machine checks the reverse half by popping matching symbols:
\begin{itemize}
    \item $1,1/\lambda$
    \item $0,0/\lambda$
\end{itemize}

\paragraph{Phase 4: accept.}
Once only the bottom-of-stack marker remains, the machine moves to $q_2$:
\begin{itemize}
    \item $\lambda,z_0/z_0$
\end{itemize}

This machine is deterministic in the sense of \sr{deterministic pushdown automaton}{deterministic pushdown automata} because the symbol $c$ tells it exactly when to stop pushing and start matching.
There is no need to guess the middle of the word.

\subsection{Accepting computation for $01c10$}

The lecture gives the following accepting sequence:
\[
(q_0,01c10,z_0)
\vdash (q_0,1c10,0z_0)
\vdash (q_0,c10,10z_0)
\vdash (q_1,10,10z_0)
\vdash (q_1,0,0z_0)
\vdash (q_1,\lambda,z_0)
\vdash (q_2,\lambda,z_0).
\]

This sequence shows the intended strategy very clearly.
The automaton first stores the prefix $01$ on the stack, then uses the marker $c$ to switch to matching mode, and finally pops the stack while reading $10$.

\subsection{Why $\{ww^R \mid w \in \{0,1\}^*\}$ is not deterministic}

The slides next compare the previous DPA with an \sr{nondeterministic pushdown automaton}{NPDA} for
\[
L = \{ww^R \mid w \in \{0,1\}^*\}.
\]

The NPDA has essentially the same push and match behaviour as before, but the move from the pushing phase to the matching phase is done using $\lambda$ rather than a visible middle marker.
That is exactly why it is nondeterministic.
The machine must guess when the first half has ended.
If it guesses too early or too late, the run fails.
Without an explicit separator such as $c$, there is no deterministic way to know where the midpoint is.

\subsection{What languages are accepted by DPAs}

Assuming \sn{final-state acceptance} for the moment, the lecture states:
\begin{itemize}
    \item every \sn{regular language} can be recognised by a DPA; and
    \item DPA languages form a strict subset of the \sr{context-free language}{context-free languages}.
\end{itemize}

The inclusion of regular languages is intuitive: a DFA can be simulated by a DPA that simply ignores its stack.
The strictness follows from languages such as $\{ww^R \mid w \in \{0,1\}^*\}$, which are context-free but not deterministic context-free.

\subsection{DPAs and unambiguous grammars}

Every language accepted by a DPA has an \sn{unambiguous grammar}, that is, a grammar in which every string has a unique \sr{Parsing?ParseTree}{parse tree}.
However, the converse is false.

The lecture gives
\[
L = \{ww^R \mid w \in \{0,1\}^*\}
\]
as an example.
This language has the unambiguous grammar
\[
S \to 0S0 \mid 1S1 \mid \lambda,
\]
but it still cannot be recognised by a DPA.
Therefore the DPA languages are a proper subset of the context-free languages that have unambiguous grammars.

\subsection{Acceptance by empty stack}

Earlier in the course, \sr{nondeterministic pushdown automaton}{NPDAs} accepting by \sn{final-state acceptance} and by \sn{empty-stack acceptance} were shown to be equivalent in expressive power.
For DPAs, this is no longer true.
The lecture emphasises that DPAs accepting by \sn{empty-stack acceptance} have very limited language-recognising power.

The key restriction is this:
\begin{quote}
If a language $L$ contains strings $x$ and $xy$ where $x \in \Sigma^*$ and $y \in \Sigma^+$, so that $x$ is a prefix of $xy$, then $L$ cannot be recognised by a DPA by empty stack.
\end{quote}

As an immediate example, the language $\{0\}^*$ cannot be recognised by a DPA by empty stack, because every accepted string is a prefix of longer accepted strings.

\subsection{Small example: $\{0,00\}$}

The slides first illustrate the argument using the language $\{0,00\}$.
Assume there is a DPA $P$ accepting this language by empty stack.
When reading the string $0$, the machine must have a run of the form
\[
(q_0,0,z_0) \vdash^* (q_1,0,\alpha) \vdash^* (q_2,\lambda,\beta) \vdash^* (q_3,\lambda,\lambda).
\]

Because the automaton is deterministic, reading the longer string $00$ must begin with exactly the same behaviour:
\[
(q_0,00,z_0) \vdash^* (q_1,00,\alpha) \vdash^* (q_2,0,\beta) \vdash^* (q_3,0,\lambda).
\]

At that point there is still one input symbol left, so the string $00$ is not yet accepted.
But the stack is already empty, so no further moves are possible.
Hence $00$ must be rejected, contradicting the assumption.

\subsection{General proof for prefix languages}

The lecture then proves the general statement.
Suppose $x \in L$ and $xy \in L$ where $y \neq \lambda$.
If a DPA $P$ accepted $L$ by empty stack, then while reading $x$ it would have to reach empty stack:
\[
(q_0,x,z_0) \vdash^* (p,\lambda,\lambda)
\]
for some state $p$.

Because the machine is deterministic, the longer input $xy$ must follow the same run prefix:
\[
(q_0,xy,z_0) \vdash^* (p,y,\lambda).
\]

Since $y \neq \lambda$, the machine has not finished reading the input.
But the stack is already empty, so no further transitions can be made.
Therefore $xy$ would be rejected, contradicting the assumption that $xy \in L$.
So no such DPA can exist.

\section{Testing Membership of a Context-Free Language}

\subsection{The problem}

The second half of the lecture asks:
given a \sr{Compilation?ContextFreeGrammar}{grammar} $G$ and a string $w$, can we test efficiently whether $w$ belongs to the language generated by $G$?

One brute-force approach would be to search through all possible \sr{Parsing?ParseTree}{parse trees} for a grammar in \sr{Chomsky normal form}{CNF} and check whether any of them derive $w$.
The lecture notes that this would be exponential in the length of $w$.

The better method is the \sr{CYK membership algorithm}{CYK algorithm}, due to Cocke, Younger, and Kasami, which runs in time $O(n^3)$.

\subsection{The basic CYK idea}

The algorithm keeps track of which non-terminals can derive which substrings of the input.
These local derivability facts are then combined until we know whether the start symbol can derive the whole string.

We assume the grammar is in \sn{Chomsky normal form}, so every \sr{Grammars?ProductionRule}{production} is of one of the forms
\[
A \to BC
\qquad \text{or} \qquad
A \to a.
\]

The lecture gives the following motivating inference.
Suppose
\[
w = w_1w_2\cdots w_n,
\]
and we already know that $N_3$ derives $w_3\cdots w_7$ and $N_8$ derives $w_8\cdots w_{11}$.
If the grammar contains the rule
\[
N_2 \to N_3N_8,
\]
then we can conclude that $N_2$ derives the longer substring $w_3\cdots w_{11}$.
In symbols:
\[
N_3 \Rightarrow^* w_3\cdots w_7,
\qquad
N_8 \Rightarrow^* w_8\cdots w_{11},
\qquad
N_2 \to N_3N_8
\]
imply
\[
N_2 \Rightarrow^* w_3\cdots w_{11}.
\]

\subsection{The CYK table}

For each substring $w_i\cdots w_j$, define $X_{ij}$ to be the set of \sr{Grammars?NonTerminalSymbol}{non-terminals} $N$ such that
\[
N \Rightarrow^* w_i\cdots w_j.
\]

These sets are arranged in a triangular table.
For a word of length $5$, the shape is:
\[
\begin{array}{ccccc}
X_{15} \\
X_{14} & X_{25} \\
X_{13} & X_{24} & X_{35} \\
X_{12} & X_{23} & X_{34} & X_{45} \\
X_{11} & X_{22} & X_{33} & X_{44} & X_{55}
\end{array}
\]
with the terminals $w_1,w_2,w_3,w_4,w_5$ below the bottom row.

Each row corresponds to substrings of a fixed length.
The bottom row has length $1$, the next row has length $2$, and so on.
The table is built from the bottom up.
The overall goal is to determine whether the start symbol $S$ belongs to the top entry $X_{1n}$.

\subsection{Base case}

If the substring length is $1$, then we are looking only at terminal productions.
So
\[
X_{ii} = \{A \mid A \to w_i \text{ is a production}\}.
\]

\subsection{Inductive step}

Now consider $X_{ij}$ for a substring of length greater than $1$.
Because the grammar is in \sr{Chomsky normal form}{CNF}, a derivation of $w_i\cdots w_j$ must begin with a rule
\[
A \to BC
\]
and a split point $k$ with $i \leq k < j$ such that
\[
B \Rightarrow^* w_i\cdots w_k
\qquad \text{and} \qquad
C \Rightarrow^* w_{k+1}\cdots w_j.
\]

Therefore $A$ belongs to $X_{ij}$ exactly when there exist $B$, $C$, and $k$ such that:
\begin{itemize}
    \item $i \leq k < j$,
    \item $B \in X_{ik}$,
    \item $C \in X_{(k+1)j}$, and
    \item $A \to BC$ is a \sr{Grammars?ProductionRule}{production rule} of the grammar.
\end{itemize}

In symbols:
\[
X_{ij}
=
\bigcup_{i \leq k < j}
\{A : A \to BC,\; B \in X_{ik},\; C \in X_{(k+1)j}\}.
\]

The slides illustrate this with examples such as:
\begin{itemize}
    \item to compute $X_{12}$, only $X_{11}$ and $X_{22}$ need to be considered;
    \item to compute $X_{24}$, we must consider both the splits $(X_{23},X_{44})$ and $(X_{22},X_{34})$.
\end{itemize}

\subsection{Complexity}

There are $n(n+1)/2$ table entries, which is $O(n^2)$.
Each entry may require checking up to $n$ split points.
Hence the overall running time is
\[
O(n^3).
\]

\section{Worked CYK Example}

The lecture works through the \sr{Compilation?ContextFreeGrammar}{grammar}
\[
S \to AB,
\qquad
A \to BB \mid a,
\qquad
B \to AB \mid b
\]
and asks whether the string $aabbb$ is derivable.

\subsection{Bottom row}

The length-$1$ substrings give:
\[
X_{11} = \{A\},
\quad
X_{22} = \{A\},
\quad
X_{33} = \{B\},
\quad
X_{44} = \{B\},
\quad
X_{55} = \{B\}.
\]

\subsection{Length-2 substrings}

Using the recurrence:
\[
X_{12} = \emptyset,
\qquad
X_{23} = \{S,B\},
\qquad
X_{34} = \{A\},
\qquad
X_{45} = \{A\}.
\]

\subsection{Length-3 substrings}

Next:
\[
X_{13} = \{S,B\},
\qquad
X_{24} = \{A\},
\qquad
X_{35} = \{S,B\}.
\]

\subsection{Length-4 substrings}

Then:
\[
X_{14} = \{A\},
\qquad
X_{25} = \{S,B\}.
\]

\subsection{Full string}

Finally:
\[
X_{15} = \{S,B\}.
\]

Since the start symbol $S$ is in $X_{15}$, the string $aabbb$ is derivable in the grammar.

\subsection{Final table}

The completed table from the slides can therefore be written as
\[
\begin{array}{ccccc}
\{S,B\} \\
\{A\} & \{S,B\} \\
\{S,B\} & \{A\} & \{S,B\} \\
\emptyset & \{S,B\} & \{A\} & \{A\} \\
\{A\} & \{A\} & \{B\} & \{B\} & \{B\}
\end{array}
\]
with terminals
\[
a \qquad a \qquad b \qquad b \qquad b
\]
underneath.

\subsection{Software note}

The lecture ends by mentioning software support.
JFLAP 7.1 includes a CYK parser, which is much faster than brute-force parsing, and JFLAP 8.0 beta includes an animation showing how the CYK table is constructed.

\section{Summary}

Lecture 17 develops three linked ideas:
\begin{itemize}
    \item the definition of \sr{deterministic pushdown automaton}{deterministic pushdown automata};
    \item the difference between \sn{final-state acceptance} and \sn{empty-stack acceptance} for DPAs; and
    \item the \sr{CYK membership algorithm}{CYK membership test} for grammars in \sn{Chomsky normal form}.
\end{itemize}

The main conceptual lesson is that determinism is a real restriction for \sr{pushdown automaton}{pushdown automata}.
It rules out some \sr{context-free language}{context-free languages} even when those languages are described by an \sn{unambiguous grammar}.
At the same time, \sr{context-free language}{context-free membership} can still be decided efficiently in \sr{Chomsky normal form}{CNF} using the \sr{CYK membership algorithm}{CYK algorithm} and its bottom-up substring table.

\end{smodule}
\end{document}
